%------------------------------------------------------------------------------%
\section{Revisão de Algumas Funções}
\label{sec:Revisao_de_Algumas_Funcoes}

Incluímos aqui uma pequena revisão de algumas funções que são importantes para o
estudo dos tópicos desse trabalho.

%------------------------------------------------------------------------------%
\subsection*{Funções Exponencial e Logarítmica}

A \emph{Função Exponencial}\index{Função!exponencial} pode ser definida de
diversas formas, por exemplo, pela sua série de Taylor
\[
  e^x = \exp(x)
      = 1 + x + \frac{x^2}{2!} + \frac{x^3}{3!} + \frac{x^4}{4!} + \cdots
      = \sum_{n=0}^{\infty} \frac{x^n}{\;n!\;}
\]
para todo $x$ real.
Consequentemente, o número $e$ é a soma da série
\[
  e = 2 + \frac{1}{2!} + \frac{1}{3!} + \frac{1}{4!} + \cdots
    = \sum_{n=0}^{\infty} \frac{1}{\;n!\;}
\]
e tem o valor
\[
 e = 2,71828\;18284\;59045\;23536\;02874\;71352\;66249\;77572\;47093\;69995\dots
\]
A imagem da função exponencial é o conjunto dos números reais positivos.
O gráfico dessa função pode ser visto na Figura~\ref{fig:functions_exp}.
Alternativamente podemos definir essa função pelo limite
\begin{equation}
  \label{eq:exponencial-pelo-limite}
  \exp(x) = \lim_{n\to\infty} \left( 1 + \frac{x}{n} \right)^n
\end{equation}
A derivada e a integral da função exponencial são
\[
 \frac{d}{dx} e^x = e^x
 \qquad\qquad
 \int e^x\,dx = e^x + C
\]

%------------------------------------------------------------------------------%
\begin{figure}[t]
  \centering
  \begin{subfigure}[b]{0.5\linewidth}
    \centering
    % Title: gl2ps_renderer figure
% Creator: GL2PS 1.4.2, (C) 1999-2020 C. Geuzaine
% For: Octave
% CreationDate: Fri Sep 10 20:09:13 2021
\setlength{\unitlength}{1pt}
\begin{picture}(0,0)
\includegraphics[width=8cm,height=4cm]{fig_functions_exp_1}
\end{picture}%
\begin{picture}(226,113)(0,0)
\fontsize{10}{0}\selectfont\put(23.3441,14.7879){\makebox(0,0)[t]{\textcolor[rgb]{0.15,0.15,0.15}{{-3}}}}
\fontsize{10}{0}\selectfont\put(53.358,14.7879){\makebox(0,0)[t]{\textcolor[rgb]{0.15,0.15,0.15}{{-2}}}}
\fontsize{10}{0}\selectfont\put(83.3719,14.7879){\makebox(0,0)[t]{\textcolor[rgb]{0.15,0.15,0.15}{{-1}}}}
\fontsize{10}{0}\selectfont\put(143.4,14.7879){\makebox(0,0)[t]{\textcolor[rgb]{0.15,0.15,0.15}{{1}}}}
\fontsize{10}{0}\selectfont\put(173.414,14.7879){\makebox(0,0)[t]{\textcolor[rgb]{0.15,0.15,0.15}{{2}}}}
\fontsize{10}{0}\selectfont\put(203.428,14.7879){\makebox(0,0)[t]{\textcolor[rgb]{0.15,0.15,0.15}{{3}}}}
\fontsize{10}{0}\selectfont\put(108.384,43.5512){\makebox(0,0)[r]{\textcolor[rgb]{0.15,0.15,0.15}{{1}}}}
\fontsize{10}{0}\selectfont\put(108.384,64.811){\makebox(0,0)[r]{\textcolor[rgb]{0.15,0.15,0.15}{{2}}}}
\fontsize{10}{0}\selectfont\put(108.384,86.0709){\makebox(0,0)[r]{\textcolor[rgb]{0.15,0.15,0.15}{{3}}}}
\fontsize{10}{0}\selectfont\put(119.389,103.079){\makebox(0,0)[l]{\textcolor[rgb]{0,0,0}{{$e^{x}$}}}}
\fontsize{10}{0}\selectfont\put(212.432,15.9134){\makebox(0,0)[l]{\textcolor[rgb]{0,0,0}{{$x$}}}}
\end{picture}

    \caption{Gráfico da função $e^{x}$}
    \label{fig:functions_exp_positiva}
  \end{subfigure}%
  \begin{subfigure}[b]{0.5\linewidth}
    \centering
    \input{fig_functions_exp_2}
    \caption{Gráfico da função $e^{-x}$}
    \label{fig:functions_exp_negativa}
  \end{subfigure}
  \caption{Gráficos da função exponencial}
  \label{fig:functions_exp}
\end{figure}
%------------------------------------------------------------------------------%

A \emph{Função Logaritmo Natural}\index{Função!logaritmo}
é a função inversa da exponencial e é o valor $y$ tal que
\[
  y = \ln(x) \qquad\Rightarrow\qquad e^y = x
\]
o domínio dessa função são os números reais positivos e sua imagem são todos os reais.
A Figura~\ref{fig:functions_log} mostra o gráfico dessa função calculada em
$\abs{x}$ estendendo assim o domínio para valores negativos de $x$.


%------------------------------------------------------------------------------%
\begin{figure}
  \centering
  % Title: gl2ps_renderer figure
% Creator: GL2PS 1.4.2, (C) 1999-2020 C. Geuzaine
% For: Octave
% CreationDate: Fri Sep 10 20:09:29 2021
\setlength{\unitlength}{1pt}
\begin{picture}(0,0)
\includegraphics[width=17cm,height=5cm]{fig_functions_log}
\end{picture}%
\begin{picture}(481,141)(0,0)
\fontsize{10}{0}\selectfont\put(87.874,62.6597){\makebox(0,0)[t]{\textcolor[rgb]{0.15,0.15,0.15}{{-6}}}}
\fontsize{10}{0}\selectfont\put(138.898,62.6597){\makebox(0,0)[t]{\textcolor[rgb]{0.15,0.15,0.15}{{-4}}}}
\fontsize{10}{0}\selectfont\put(189.921,62.6597){\makebox(0,0)[t]{\textcolor[rgb]{0.15,0.15,0.15}{{-2}}}}
\fontsize{10}{0}\selectfont\put(291.969,62.6597){\makebox(0,0)[t]{\textcolor[rgb]{0.15,0.15,0.15}{{2}}}}
\fontsize{10}{0}\selectfont\put(342.992,62.6597){\makebox(0,0)[t]{\textcolor[rgb]{0.15,0.15,0.15}{{4}}}}
\fontsize{10}{0}\selectfont\put(394.016,62.6597){\makebox(0,0)[t]{\textcolor[rgb]{0.15,0.15,0.15}{{6}}}}
\fontsize{10}{0}\selectfont\put(235.944,19.1102){\makebox(0,0)[r]{\textcolor[rgb]{0.15,0.15,0.15}{{-2}}}}
\fontsize{10}{0}\selectfont\put(235.944,44.6221){\makebox(0,0)[r]{\textcolor[rgb]{0.15,0.15,0.15}{{-1}}}}
\fontsize{10}{0}\selectfont\put(235.944,95.6457){\makebox(0,0)[r]{\textcolor[rgb]{0.15,0.15,0.15}{{1}}}}
\fontsize{10}{0}\selectfont\put(235.944,121.157){\makebox(0,0)[r]{\textcolor[rgb]{0.15,0.15,0.15}{{2}}}}
\fontsize{10}{0}\selectfont\put(246.047,126.26){\makebox(0,0)[l]{\textcolor[rgb]{0,0,0}{{$\ln\abs{x}$}}}}
\fontsize{10}{0}\selectfont\put(414.425,62.4803){\makebox(0,0)[l]{\textcolor[rgb]{0,0,0}{{$x$}}}}
\end{picture}

  \caption{Gráfico da função $\ln|x|$}
  \label{fig:functions_log}
\end{figure}
%------------------------------------------------------------------------------%

A série de Taylor da função logaritmo normalmente é escrita como
\[
  \ln(x+1)
  = x - \frac{x^2}{2}
  + \frac{x^3}{3}
  - \frac{x^4}{4}
  + \cdots
  = \sum_{n=1}^{\infty} \frac{(-1)^{n-1}x^n}{n}
\]
e tem raio de convergência $R=1$.
A derivada e a integral da função exponencial são
\[
  \frac{d}{dx} \ln(x) = \frac{1}{\,x\,}
  \qquad\qquad
  \int \ln(x)\,dx = x\big( \ln(x) - 1 \big) + C
\]

%------------------------------------------------------------------------------%
\subsection*{Funções Trigonométricas}

Tradicionalmente pensamos nas funções trigonométricas como combinações
das funções seno e cosseno
\begin{align*}
	\tan(x) & = \frac{\sin(x)}{\cos(x)} &
	\cot(x) & = \frac{\cos(x)}{\sin(x)} = \frac{1}{\tan(x)}\\[5mm]
	\sec(x) & = \frac{1}{\cos(x)}       &
	\csc(x) & = \frac{1}{\sin(x)}
\end{align*}
Os gráficos dessas funções podem ser vistos na Figura~\ref{fig:functions_trig}.

%------------------------------------------------------------------------------%
\begin{figure}
  \centering
  %--------------------------------------------------------------%
  \begin{subfigure}[b]{0.5\linewidth}
    \centering
    \input{fig_functions_sin}
    \caption{Gráfico da função $\sin(x)$}
    \label{fig:functions_sin}
  \end{subfigure}%
  \begin{subfigure}[b]{0.5\linewidth}
    \centering
    % Title: gl2ps_renderer figure
% Creator: GL2PS 1.4.2, (C) 1999-2020 C. Geuzaine
% For: Octave
% CreationDate: Fri Sep 10 20:09:21 2021
\setlength{\unitlength}{1pt}
\begin{picture}(0,0)
\includegraphics[width=8cm,height=4cm]{fig_functions_cos}
\end{picture}%
\begin{picture}(226,113)(0,0)
\fontsize{10}{0}\selectfont\put(23.3441,48.8036){\makebox(0,0)[t]{\textcolor[rgb]{0.15,0.15,0.15}{{$-3 \pi$}}}}
\fontsize{10}{0}\selectfont\put(53.358,48.8036){\makebox(0,0)[t]{\textcolor[rgb]{0.15,0.15,0.15}{{$-2 \pi$}}}}
\fontsize{10}{0}\selectfont\put(83.3719,48.8036){\makebox(0,0)[t]{\textcolor[rgb]{0.15,0.15,0.15}{{$-\pi$}}}}
\fontsize{10}{0}\selectfont\put(143.4,48.8036){\makebox(0,0)[t]{\textcolor[rgb]{0.15,0.15,0.15}{{$\pi$}}}}
\fontsize{10}{0}\selectfont\put(173.414,48.8036){\makebox(0,0)[t]{\textcolor[rgb]{0.15,0.15,0.15}{{$2 \pi$}}}}
\fontsize{10}{0}\selectfont\put(203.428,48.8036){\makebox(0,0)[t]{\textcolor[rgb]{0.15,0.15,0.15}{{$3 \pi$}}}}
\fontsize{10}{0}\selectfont\put(108.384,26.2932){\makebox(0,0)[r]{\textcolor[rgb]{0.15,0.15,0.15}{{-1}}}}
\fontsize{10}{0}\selectfont\put(108.384,86.321){\makebox(0,0)[r]{\textcolor[rgb]{0.15,0.15,0.15}{{1}}}}
\fontsize{10}{0}\selectfont\put(119.118,99.8272){\makebox(0,0)[l]{\textcolor[rgb]{0,0,0}{{$\cos(x)$}}}}
\fontsize{10}{0}\selectfont\put(212.432,50.3043){\makebox(0,0)[l]{\textcolor[rgb]{0,0,0}{{$x$}}}}
\end{picture}

    \caption{Gráfico da função $\cos(x)$}
    \label{fig:functions_cos}
  \end{subfigure}\\[0.5\baselineskip]
  %--------------------------------------------------------------%
  \begin{subfigure}[b]{0.5\linewidth}
    \centering
    % Title: gl2ps_renderer figure
% Creator: GL2PS 1.4.2, (C) 1999-2020 C. Geuzaine
% For: Octave
% CreationDate: Fri Sep 10 20:09:21 2021
\setlength{\unitlength}{1pt}
\begin{picture}(0,0)
\includegraphics[width=8cm,height=4cm]{fig_functions_tan}
\end{picture}%
\begin{picture}(226,113)(0,0)
\fontsize{10}{0}\selectfont\put(45.3543,48.8036){\makebox(0,0)[t]{\textcolor[rgb]{0.15,0.15,0.15}{{$-\pi$}}}}
\fontsize{10}{0}\selectfont\put(181.417,48.8036){\makebox(0,0)[t]{\textcolor[rgb]{0.15,0.15,0.15}{{$\pi$}}}}
\fontsize{10}{0}\selectfont\put(108.384,6.52794){\makebox(0,0)[r]{\textcolor[rgb]{0.15,0.15,0.15}{{-4}}}}
\fontsize{10}{0}\selectfont\put(108.384,43.8623){\makebox(0,0)[r]{\textcolor[rgb]{0.15,0.15,0.15}{{-1}}}}
\fontsize{10}{0}\selectfont\put(108.384,68.7519){\makebox(0,0)[r]{\textcolor[rgb]{0.15,0.15,0.15}{{1}}}}
\fontsize{10}{0}\selectfont\put(108.384,106.086){\makebox(0,0)[r]{\textcolor[rgb]{0.15,0.15,0.15}{{4}}}}
\fontsize{10}{0}\selectfont\put(161.027,106.086){\makebox(0,0)[l]{\textcolor[rgb]{0,0,0}{{$\tan(x)$}}}}
\fontsize{10}{0}\selectfont\put(212.031,50.0847){\makebox(0,0)[l]{\textcolor[rgb]{0,0,0}{{$x$}}}}
\end{picture}

    \caption{Gráfico da função $\tan(x)$}
    \label{fig:functions_tan}
  \end{subfigure}%
  \begin{subfigure}[b]{0.5\linewidth}
    \centering
    \input{fig_functions_cot}
    \caption{Gráfico da função $\cot(x)$}
    \label{fig:functions_cot}
  \end{subfigure}\\[0.5\baselineskip]
  %--------------------------------------------------------------%
  \begin{subfigure}[b]{0.5\linewidth}
    \centering
    \input{fig_functions_sec}
    \caption{Gráfico da função $\sec(x)$}
    \label{fig:functions_sec}
  \end{subfigure}%
  \begin{subfigure}[b]{0.5\linewidth}
    \centering
    % Title: gl2ps_renderer figure
% Creator: GL2PS 1.4.2, (C) 1999-2020 C. Geuzaine
% For: Octave
% CreationDate: Fri Sep 10 20:09:23 2021
\setlength{\unitlength}{1pt}
\begin{picture}(0,0)
\includegraphics[width=8cm,height=4cm]{fig_functions_csc}
\end{picture}%
\begin{picture}(226,113)(0,0)
\fontsize{10}{0}\selectfont\put(20.6156,48.8036){\makebox(0,0)[t]{\textcolor[rgb]{0.15,0.15,0.15}{{$-2 \pi$}}}}
\fontsize{10}{0}\selectfont\put(67.0007,48.8036){\makebox(0,0)[t]{\textcolor[rgb]{0.15,0.15,0.15}{{$-\pi$}}}}
\fontsize{10}{0}\selectfont\put(159.771,48.8036){\makebox(0,0)[t]{\textcolor[rgb]{0.15,0.15,0.15}{{$\pi$}}}}
\fontsize{10}{0}\selectfont\put(206.156,48.8036){\makebox(0,0)[t]{\textcolor[rgb]{0.15,0.15,0.15}{{$2 \pi$}}}}
\fontsize{10}{0}\selectfont\put(108.384,6.52794){\makebox(0,0)[r]{\textcolor[rgb]{0.15,0.15,0.15}{{-4}}}}
\fontsize{10}{0}\selectfont\put(108.384,43.8623){\makebox(0,0)[r]{\textcolor[rgb]{0.15,0.15,0.15}{{-1}}}}
\fontsize{10}{0}\selectfont\put(108.384,68.7519){\makebox(0,0)[r]{\textcolor[rgb]{0.15,0.15,0.15}{{1}}}}
\fontsize{10}{0}\selectfont\put(108.384,106.086){\makebox(0,0)[r]{\textcolor[rgb]{0.15,0.15,0.15}{{4}}}}
\fontsize{10}{0}\selectfont\put(163.586,106.086){\makebox(0,0)[l]{\textcolor[rgb]{0,0,0}{{$\csc(x)$}}}}
\fontsize{10}{0}\selectfont\put(215.433,52.5737){\makebox(0,0)[l]{\textcolor[rgb]{0,0,0}{{$x$}}}}
\end{picture}

    \caption{Gráfico da função $\csc(x)$}
    \label{fig:functions_csc}
  \end{subfigure}\\[0.5\baselineskip]
  %--------------------------------------------------------------%
  \caption{Gráficos das funções trigonométricas}
  \label{fig:functions_trig}
\end{figure}
%------------------------------------------------------------------------------%

Identidade Pitagórica
\begin{align*}
  \sin^2(x) + \cos^2(x) & = 1         &
  1 + \tan^2(x)         & = \sec^2(x) &
  1 + \cot^2(x)         & = \csc^2(x)
\end{align*}
Algumas identidades trigonométricas
\begin{align*}
	& \sin(-x)                         = -\sin(x)  &
	& \cos(-x)                         = \cos(x)   &
	& \tan(-x)                         = -\tan(x)  \\[5mm]
	& \sin\left(\frac{\pi}{2}-x\right) = \cos(x)   &
	& \cos\left(\frac{\pi}{2}-x\right) = \sin(x)   &
	& \tan\left(\frac{\pi}{2}-x\right) = \cot(x)
\end{align*}
Fórmulas de soma e subtração
\begin{align*}
	\sin(x+y) & = \sin(x)\cos(y) + \sin(y)\cos(x) \\
	\sin(x-y) & = \sin(x)\cos(y) - \sin(y)\cos(x) \\
	\cos(x+y) & = \cos(x)\cos(y) - \sin(x)\sin(y) \\
	\cos(x-y) & = \cos(x)\cos(y) + \sin(x)\sin(y)
\end{align*}
Fórmulas de ângulo duplo
\begin{align*}
	\cos(2x) & = \cos^2(x) - \sin^2(x) & 	\sin(2x) & = 2\sin(x)\cos(x)  \\
	         & = 2\cos^2(x) - 1        \\
	         & = 1 - 2\sin^2(x)
\end{align*}
Fórmulas de metade do ângulo
\begin{align*}
	\sin^2(\frac{x}{2}) & = \frac{1-\cos(x)}{2} &
  \cos^2(\frac{x}{2}) & = \frac{1+\cos(x)}{2}
\end{align*}

%------------------------------------------------------------------------------%
\begin{figure}
  \centering
  %--------------------------------------------------------------%
  \begin{subfigure}[b]{0.5\linewidth}
    \centering
    \input{fig_functions_asin}
    \caption{Gráfico da função $\arcsin(x)$}
    \label{fig:functions_asin}
  \end{subfigure}%
  \begin{subfigure}[b]{0.5\linewidth}
    \centering
    % Title: gl2ps_renderer figure
% Creator: GL2PS 1.4.2, (C) 1999-2020 C. Geuzaine
% For: Octave
% CreationDate: Fri Sep 10 20:09:16 2021
\setlength{\unitlength}{1pt}
\begin{picture}(0,0)
\includegraphics[width=8cm,height=4cm]{fig_functions_acos}
\end{picture}%
\begin{picture}(226,113)(0,0)
\fontsize{10}{0}\selectfont\put(20.6156,3.36282){\makebox(0,0)[t]{\textcolor[rgb]{0.15,0.15,0.15}{{-1}}}}
\fontsize{10}{0}\selectfont\put(67.0007,3.36282){\makebox(0,0)[t]{\textcolor[rgb]{0.15,0.15,0.15}{{-0,5}}}}
\fontsize{10}{0}\selectfont\put(159.771,3.36282){\makebox(0,0)[t]{\textcolor[rgb]{0.15,0.15,0.15}{{0,5}}}}
\fontsize{10}{0}\selectfont\put(206.156,3.36282){\makebox(0,0)[t]{\textcolor[rgb]{0.15,0.15,0.15}{{1}}}}
\fontsize{10}{0}\selectfont\put(108.384,54.7138){\makebox(0,0)[r]{\textcolor[rgb]{0.15,0.15,0.15}{{$ \sfrac{   \pi}{2}$}}}}
\fontsize{10}{0}\selectfont\put(108.384,98.5612){\makebox(0,0)[r]{\textcolor[rgb]{0.15,0.15,0.15}{{$\pi$}}}}
\fontsize{10}{0}\selectfont\put(122.663,102.983){\makebox(0,0)[l]{\textcolor[rgb]{0,0,0}{{$\arccos(x)$}}}}
\fontsize{10}{0}\selectfont\put(215.433,2.49205){\makebox(0,0)[l]{\textcolor[rgb]{0,0,0}{{$x$}}}}
\end{picture}

    \caption{Gráfico da função $\arccos(x)$}
    \label{fig:functions_acos}
  \end{subfigure}\\[0.5\baselineskip]
  %--------------------------------------------------------------%
  \begin{subfigure}[b]{0.5\linewidth}
    \centering
    % Title: gl2ps_renderer figure
% Creator: GL2PS 1.4.2, (C) 1999-2020 C. Geuzaine
% For: Octave
% CreationDate: Fri Sep 10 20:09:16 2021
\setlength{\unitlength}{1pt}
\begin{picture}(0,0)
\includegraphics[width=8cm,height=4cm]{fig_functions_atan}
\end{picture}%
\begin{picture}(226,113)(0,0)
\fontsize{10}{0}\selectfont\put(20.6156,48.8036){\makebox(0,0)[t]{\textcolor[rgb]{0.15,0.15,0.15}{{-10}}}}
\fontsize{10}{0}\selectfont\put(67.0007,48.8036){\makebox(0,0)[t]{\textcolor[rgb]{0.15,0.15,0.15}{{-5}}}}
\fontsize{10}{0}\selectfont\put(159.771,48.8036){\makebox(0,0)[t]{\textcolor[rgb]{0.15,0.15,0.15}{{5}}}}
\fontsize{10}{0}\selectfont\put(206.156,48.8036){\makebox(0,0)[t]{\textcolor[rgb]{0.15,0.15,0.15}{{10}}}}
\fontsize{10}{0}\selectfont\put(108.384,13.7874){\makebox(0,0)[r]{\textcolor[rgb]{0.15,0.15,0.15}{{$-\sfrac{   \pi}{2}$}}}}
\fontsize{10}{0}\selectfont\put(108.384,98.8268){\makebox(0,0)[r]{\textcolor[rgb]{0.15,0.15,0.15}{{$ \sfrac{   \pi}{2}$}}}}
\fontsize{10}{0}\selectfont\put(118.024,107.331){\makebox(0,0)[l]{\textcolor[rgb]{0,0,0}{{$\arctan(x)$}}}}
\fontsize{10}{0}\selectfont\put(215.433,48.1864){\makebox(0,0)[l]{\textcolor[rgb]{0,0,0}{{$x$}}}}
\end{picture}

    \caption{Gráfico da função $\arctan(x)$}
    \label{fig:functions_atan}
  \end{subfigure}%
  \begin{subfigure}[b]{0.5\linewidth}
    \centering
    % Title: gl2ps_renderer figure
% Creator: GL2PS 1.4.2, (C) 1999-2020 C. Geuzaine
% For: Octave
% CreationDate: Fri Sep 10 20:09:16 2021
\setlength{\unitlength}{1pt}
\begin{picture}(0,0)
\includegraphics[width=8cm,height=4cm]{fig_functions_acot}
\end{picture}%
\begin{picture}(226,113)(0,0)
\fontsize{10}{0}\selectfont\put(20.6156,3.36282){\makebox(0,0)[t]{\textcolor[rgb]{0.15,0.15,0.15}{{-10}}}}
\fontsize{10}{0}\selectfont\put(67.0007,3.36282){\makebox(0,0)[t]{\textcolor[rgb]{0.15,0.15,0.15}{{-5}}}}
\fontsize{10}{0}\selectfont\put(159.771,3.36282){\makebox(0,0)[t]{\textcolor[rgb]{0.15,0.15,0.15}{{5}}}}
\fontsize{10}{0}\selectfont\put(206.156,3.36282){\makebox(0,0)[t]{\textcolor[rgb]{0.15,0.15,0.15}{{10}}}}
\fontsize{10}{0}\selectfont\put(108.384,54.7138){\makebox(0,0)[r]{\textcolor[rgb]{0.15,0.15,0.15}{{$ \sfrac{   \pi}{2}$}}}}
\fontsize{10}{0}\selectfont\put(108.384,98.5612){\makebox(0,0)[r]{\textcolor[rgb]{0.15,0.15,0.15}{{$\pi$}}}}
\fontsize{10}{0}\selectfont\put(118.024,107.331){\makebox(0,0)[l]{\textcolor[rgb]{0,0,0}{{$\arccot(x)$}}}}
\fontsize{10}{0}\selectfont\put(215.433,2.49205){\makebox(0,0)[l]{\textcolor[rgb]{0,0,0}{{$x$}}}}
\end{picture}

    \caption{Gráfico da função $\arccot(x)$}
    \label{fig:functions_acot}
  \end{subfigure}\\[0.5\baselineskip]
  %--------------------------------------------------------------%
  \begin{subfigure}[b]{0.5\linewidth}
    \centering
    % Title: gl2ps_renderer figure
% Creator: GL2PS 1.4.2, (C) 1999-2020 C. Geuzaine
% For: Octave
% CreationDate: Fri Sep 10 20:09:17 2021
\setlength{\unitlength}{1pt}
\begin{picture}(0,0)
\includegraphics[width=8cm,height=4cm]{fig_functions_asec}
\end{picture}%
\begin{picture}(226,113)(0,0)
\fontsize{10}{0}\selectfont\put(20.6156,7.5689){\makebox(0,0)[t]{\textcolor[rgb]{0.15,0.15,0.15}{{-10}}}}
\fontsize{10}{0}\selectfont\put(67.0007,7.5689){\makebox(0,0)[t]{\textcolor[rgb]{0.15,0.15,0.15}{{-5}}}}
\fontsize{10}{0}\selectfont\put(104.109,7.5689){\makebox(0,0)[t]{\textcolor[rgb]{0.15,0.15,0.15}{{-1}}}}
\fontsize{10}{0}\selectfont\put(122.663,7.5689){\makebox(0,0)[t]{\textcolor[rgb]{0.15,0.15,0.15}{{1}}}}
\fontsize{10}{0}\selectfont\put(159.771,7.5689){\makebox(0,0)[t]{\textcolor[rgb]{0.15,0.15,0.15}{{5}}}}
\fontsize{10}{0}\selectfont\put(206.156,7.5689){\makebox(0,0)[t]{\textcolor[rgb]{0.15,0.15,0.15}{{10}}}}
\fontsize{10}{0}\selectfont\put(108.384,91.9543){\makebox(0,0)[r]{\textcolor[rgb]{0.15,0.15,0.15}{{$ \sfrac{   \pi}{2}$}}}}
\fontsize{10}{0}\selectfont\put(118.024,107.331){\makebox(0,0)[l]{\textcolor[rgb]{0,0,0}{{$\arcsec(x)$}}}}
\fontsize{10}{0}\selectfont\put(215.433,5.28347){\makebox(0,0)[l]{\textcolor[rgb]{0,0,0}{{$x$}}}}
\end{picture}

    \caption{Gráfico da função $\arcsec(x)$}
    \label{fig:functions_asec}
  \end{subfigure}%
  \begin{subfigure}[b]{0.5\linewidth}
    \centering
    \input{fig_functions_acsc}
    \caption{Gráfico da função $\arccsc(x)$}
    \label{fig:functions_acsc}
  \end{subfigure}\\[0.5\baselineskip]
  %--------------------------------------------------------------%
  \caption{Gráficos das funções trigonométricas inversas}
  \label{fig:functions_trig_inv}
\end{figure}
%------------------------------------------------------------------------------%

%------------------------------------------------------------------------------%
\subsection*{Funções Trigonométricas Inversas}

Os gráficos dessas funções podem ser vistos na Figura~\ref{fig:functions_trig_inv}.

Derivadas as funções trigonométricas inversas
\begin{align*}
  \frac{d}{dx} \arcsin(x) & =  \frac{1}{\sqrt{1-x^2\,}\,}  &
  \frac{d}{dx} \arccos(x) & = -\frac{1}{\sqrt{1-x^2\,}\,}  \\[3mm]
  \frac{d}{dx} \arctan(x) & =  \frac{1}{\,1-x^2\,}         &
  \frac{d}{dx} \arccot(x) & = -\frac{1}{\,1-x^2\,}         \\[3mm]
  \frac{d}{dx} \arcsec(x) & =  \frac{1}{x\sqrt{x^2-1\,}\,} &
  \frac{d}{dx} \arccsc(x) & = -\frac{1}{x\sqrt{x^2-1\,}\,}
\end{align*}

\FloatBarrier

%------------------------------------------------------------------------------%
\subsection*{Funções Hiperbólicas}

As \emph{funções hiperbólicas}\index{Função!hiperbólica}
são funções análogas às funções trigonométricas ordinárias,
que recebem esse nome por estarem relacionadas às hipérboles de modo similar
a relação entre as funções trigonométrica e a circunferência. Porém, em nosso
contexto essa interpretação geométrica não será utilizada. Aqui elas são importantes
por representarem funções exponenciais em um arranjo similar às trigonométricas.

Definição do seno e cosseno hiperbólicos%
\index{Seno hiperbólico}
\index{Cosseno hiperbólico}
\begin{align}
\sinh(x) & = \frac{e^x - e^{-x}}{2} \label{eq:def_seno_hiperbolico}    \\[3mm]
\cosh(x) & = \frac{e^x + e^{-x}}{2} \label{eq:def_cosseno_hiperbolico}
\end{align}
Os gráficos dessas funções podem ser vistos na Figura~\ref{fig:functions_hyper}.

%------------------------------------------------------------------------------%
\begin{figure}
  \centering
  \begin{subfigure}[b]{0.5\linewidth}
    \centering
    % Title: gl2ps_renderer figure
% Creator: GL2PS 1.4.2, (C) 1999-2020 C. Geuzaine
% For: Octave
% CreationDate: Fri Sep 10 20:09:26 2021
\setlength{\unitlength}{1pt}
\begin{picture}(0,0)
\includegraphics[width=8cm,height=8cm]{fig_functions_sinh}
\end{picture}%
\begin{picture}(226,226)(0,0)
\fontsize{10}{0}\selectfont\put(45.3543,105.111){\makebox(0,0)[t]{\textcolor[rgb]{0.15,0.15,0.15}{{-2}}}}
\fontsize{10}{0}\selectfont\put(79.3701,105.111){\makebox(0,0)[t]{\textcolor[rgb]{0.15,0.15,0.15}{{-1}}}}
\fontsize{10}{0}\selectfont\put(147.402,105.111){\makebox(0,0)[t]{\textcolor[rgb]{0.15,0.15,0.15}{{1}}}}
\fontsize{10}{0}\selectfont\put(181.417,105.111){\makebox(0,0)[t]{\textcolor[rgb]{0.15,0.15,0.15}{{2}}}}
\fontsize{10}{0}\selectfont\put(108.384,19.8439){\makebox(0,0)[r]{\textcolor[rgb]{0.15,0.15,0.15}{{-10}}}}
\fontsize{10}{0}\selectfont\put(108.384,66.2291){\makebox(0,0)[r]{\textcolor[rgb]{0.15,0.15,0.15}{{-5}}}}
\fontsize{10}{0}\selectfont\put(108.384,158.999){\makebox(0,0)[r]{\textcolor[rgb]{0.15,0.15,0.15}{{5}}}}
\fontsize{10}{0}\selectfont\put(108.384,205.384){\makebox(0,0)[r]{\textcolor[rgb]{0.15,0.15,0.15}{{10}}}}
\fontsize{10}{0}\selectfont\put(123.591,210.023){\makebox(0,0)[l]{\textcolor[rgb]{0,0,0}{{$\sinh(x)$}}}}
\fontsize{10}{0}\selectfont\put(208.63,105.193){\makebox(0,0)[l]{\textcolor[rgb]{0,0,0}{{$x$}}}}
\end{picture}

    \caption{Gráfico da função $\sinh(x)$}
    \label{fig:functions_sinh}
  \end{subfigure}%
  \begin{subfigure}[b]{0.5\linewidth}
    \centering
    \input{fig_functions_cosh}
    \caption{Gráfico da função $\cosh(x)$}
    \label{fig:functions_cosh}
  \end{subfigure}
  \caption{Gráficos das funções hiperbólicas}
  \label{fig:functions_hyper}
\end{figure}
%------------------------------------------------------------------------------%

Também podemos definir essas funções com o uso de expressões envolvendo números complexos
\begin{align*}
\sinh(x) & = -i \sin(ix) &
\cosh(x) & =    \cos(ix)
\end{align*}

Derivadas das funções hiperbólicas
\begin{align*}
\frac{d  }{dx  } \sinh(x) & = \cosh(x) & \frac{d  }{dx  } \cosh(x) & = \sinh(x) \\[3mm]
\frac{d^2}{dx^2} \sinh(x) & = \sinh(x) & \frac{d^2}{dx^2} \cosh(x) & = \cosh(x)
\end{align*}
essas derivadas podem ser calculadas de modo simples apenas substituindo a
definição da função, por exemplo
\begin{align*}
	\frac{d}{dx} \sinh(x) & = \frac{d}{dx} \frac{e^x - e^{-x}}{2}                             \\[3mm]
	                      & = \frac{1}{2} \left( \frac{de^x}{dx} - \frac{de^{-x}}{dx} \right) \\[3mm]
	                      & = \frac{1}{2} \left( e^x + e^{-x} \right)                         \\[3mm]
	                      & = \cosh(x)
\end{align*}
Algumas relações úteis
\begin{align*}
\sinh(-x) & = -\sinh(x) &
\cosh(-x) & =  \cosh(x)
\end{align*}
Relação com a exponencial e \emph{Fórmula de Euler}\index{Fórmula!Euler}
\begin{align*}
	e^x     & = \cosh(x) + \sinh(x) \\
	e^{-x}  & = \cosh(x) - \sinh(x) \\
	e^{-ix} & = \cos(x) + i\sin(x)
\end{align*}
Identidade similar a relação pitagórica
\[ \cosh^2(x) - \sinh^1(x) = 1 \]
Somando e subtraindo os argumentos
\begin{gather*}
\sinh(x+y) = \sinh(x)\cosh(y) + \cosh(x)\sinh(y) \\
\sinh(x-y) = \sinh(x)\cosh(y) - \cosh(x)\sinh(y) \\
\cosh(x+y) = \cosh(x)\cosh(y) + \sinh(x)\sinh(y) \\
\cosh(x-y) = \cosh(x)\cosh(y) - \sinh(x)\sinh(y)
\end{gather*}
Séries de Taylor
\begin{gather*}
\sinh(x) = x + \frac{x^3}{3!} + \frac{x^5}{5!} + \frac{x^7}{7!} + \cdots
         = \sum_{n=0}^{\infty} \frac{x^{2n+1}}{\;(2n+1)!\;} \\
\cosh(x) = 1 + \frac{x^2}{2!} + \frac{x^4}{4!} + \frac{x^6}{6!} + \cdots
         = \sum_{n=0}^{\infty} \frac{x^{2n}}{\;(2n)!\;}
\end{gather*}
Da mesma forma como as funções trigonométricas podem ser definidas a partir do seno e
cosseno, podemos definir as seguintes funções hiperbólicas
\begin{align*}
  \tanh(x) & = \frac{\sinh(x)}{\cosh(x)} &
  \coth(x) & = \frac{\cosh(x)}{\sinh(x)} = \frac{1}{\tanh(x)}\\[5mm]
  \sech(x) & = \frac{1}{\cosh(x)}       &
  \csch(x) & = \frac{1}{\sinh(x)}
\end{align*}
Os gráficos dessas funções podem ser vistos na Figura~\ref{fig:functions_hiper_2}.

%------------------------------------------------------------------------------%
\begin{figure}
  \centering
  %--------------------------------------------------------------%
  \begin{subfigure}[b]{0.5\linewidth}
    \centering
    \input{fig_functions_tanh}
    \caption{Gráfico da função $\tanh(x)$}
    \label{fig:functions_tanh}
  \end{subfigure}%
  \begin{subfigure}[b]{0.5\linewidth}
    \centering
    \input{fig_functions_coth}
    \caption{Gráfico da função $\coth(x)$}
    \label{fig:functions_coth}
  \end{subfigure}\\[0.5\baselineskip]
  %--------------------------------------------------------------%
  \begin{subfigure}[b]{0.5\linewidth}
    \centering
    % Title: gl2ps_renderer figure
% Creator: GL2PS 1.4.2, (C) 1999-2020 C. Geuzaine
% For: Octave
% CreationDate: Tue Oct 26 12:50:34 2021
\setlength{\unitlength}{1pt}
\begin{picture}(0,0)
\includegraphics[width=8cm,height=4cm]{fig_functions_sech}
\end{picture}%
\begin{picture}(226,113)(0,0)
\fontsize{10}{0}\selectfont\put(45.3543,48.8036){\makebox(0,0)[t]{\textcolor[rgb]{0.15,0.15,0.15}{{-2}}}}
\fontsize{10}{0}\selectfont\put(79.3701,48.8036){\makebox(0,0)[t]{\textcolor[rgb]{0.15,0.15,0.15}{{-1}}}}
\fontsize{10}{0}\selectfont\put(147.402,48.8036){\makebox(0,0)[t]{\textcolor[rgb]{0.15,0.15,0.15}{{1}}}}
\fontsize{10}{0}\selectfont\put(181.417,48.8036){\makebox(0,0)[t]{\textcolor[rgb]{0.15,0.15,0.15}{{2}}}}
\fontsize{10}{0}\selectfont\put(108.384,17.0581){\makebox(0,0)[r]{\textcolor[rgb]{0.15,0.15,0.15}{{-1}}}}
\fontsize{10}{0}\selectfont\put(108.384,95.556){\makebox(0,0)[r]{\textcolor[rgb]{0.15,0.15,0.15}{{1}}}}
\fontsize{10}{0}\selectfont\put(123.591,107.331){\makebox(0,0)[l]{\textcolor[rgb]{0,0,0}{{$\sech(x)$}}}}
\fontsize{10}{0}\selectfont\put(208.63,52.3822){\makebox(0,0)[l]{\textcolor[rgb]{0,0,0}{{$x$}}}}
\end{picture}

    \caption{Gráfico da função $\sech(x)$}
    \label{fig:functions_sech}
  \end{subfigure}%
  \begin{subfigure}[b]{0.5\linewidth}
    \centering
    % Title: gl2ps_renderer figure
% Creator: GL2PS 1.4.2, (C) 1999-2020 C. Geuzaine
% For: Octave
% CreationDate: Tue Oct 26 12:50:35 2021
\setlength{\unitlength}{1pt}
\begin{picture}(0,0)
\includegraphics[width=8cm,height=4cm]{fig_functions_csch}
\end{picture}%
\begin{picture}(226,113)(0,0)
\fontsize{10}{0}\selectfont\put(45.3543,48.8036){\makebox(0,0)[t]{\textcolor[rgb]{0.15,0.15,0.15}{{-2}}}}
\fontsize{10}{0}\selectfont\put(79.3701,48.8036){\makebox(0,0)[t]{\textcolor[rgb]{0.15,0.15,0.15}{{-1}}}}
\fontsize{10}{0}\selectfont\put(147.402,48.8036){\makebox(0,0)[t]{\textcolor[rgb]{0.15,0.15,0.15}{{1}}}}
\fontsize{10}{0}\selectfont\put(181.417,48.8036){\makebox(0,0)[t]{\textcolor[rgb]{0.15,0.15,0.15}{{2}}}}
\fontsize{10}{0}\selectfont\put(108.384,6.52794){\makebox(0,0)[r]{\textcolor[rgb]{0.15,0.15,0.15}{{-4}}}}
\fontsize{10}{0}\selectfont\put(108.384,43.8623){\makebox(0,0)[r]{\textcolor[rgb]{0.15,0.15,0.15}{{-1}}}}
\fontsize{10}{0}\selectfont\put(108.384,68.7519){\makebox(0,0)[r]{\textcolor[rgb]{0.15,0.15,0.15}{{1}}}}
\fontsize{10}{0}\selectfont\put(108.384,106.086){\makebox(0,0)[r]{\textcolor[rgb]{0.15,0.15,0.15}{{4}}}}
\fontsize{10}{0}\selectfont\put(126.992,106.086){\makebox(0,0)[l]{\textcolor[rgb]{0,0,0}{{$\csch(x)$}}}}
\fontsize{10}{0}\selectfont\put(208.63,52.5737){\makebox(0,0)[l]{\textcolor[rgb]{0,0,0}{{$x$}}}}
\end{picture}

    \caption{Gráfico da função $\csch(x)$}
    \label{fig:functions_csch}
  \end{subfigure}\\[0.5\baselineskip]
  %--------------------------------------------------------------%
  \caption{Gráficos das funções hiperbólicas compostas}
  \label{fig:functions_hiper_2}
\end{figure}
%------------------------------------------------------------------------------%

%------------------------------------------------------------------------------%
\subsection*{Função sinc}

O termo \emph{sinc}\index{sinc} é uma contração do nome da função em latim
\textit{sinus cardinalis} (seno cardinal), essa função é definida como
\begin{equation}
  \label{eq:funcao_sinc}
  \sinc(x) = \begin{cases}
                \dfrac{\sin(x)}{x} & \; x \neq 0 \\[3mm]
                \quad 1            & \; x = 0
             \end{cases}
\end{equation}
Porém, por simplicidade, é comum escrevermos apenas
\[
  \sinc(x) = \frac{\sin(x)}{x}
\]
O gráfico dessa função pode ser visto na Figura~\ref{fig:functions_sinc}.

%------------------------------------------------------------------------------%
\begin{figure}
  \centering
  % Title: gl2ps_renderer figure
% Creator: GL2PS 1.4.2, (C) 1999-2020 C. Geuzaine
% For: Octave
% CreationDate: Fri Sep 10 20:09:31 2021
\setlength{\unitlength}{1pt}
\begin{picture}(0,0)
\includegraphics[width=17cm,height=6cm]{fig_functions_sinc}
\end{picture}%
\begin{picture}(481,170)(0,0)
\fontsize{10}{0}\selectfont\put(49.6063,33.7226){\makebox(0,0)[t]{\textcolor[rgb]{0.15,0.15,0.15}{{$-6 \pi$}}}}
\fontsize{10}{0}\selectfont\put(81.4961,33.7226){\makebox(0,0)[t]{\textcolor[rgb]{0.15,0.15,0.15}{{$-5 \pi$}}}}
\fontsize{10}{0}\selectfont\put(113.386,33.7226){\makebox(0,0)[t]{\textcolor[rgb]{0.15,0.15,0.15}{{$-4 \pi$}}}}
\fontsize{10}{0}\selectfont\put(145.276,33.7226){\makebox(0,0)[t]{\textcolor[rgb]{0.15,0.15,0.15}{{$-3 \pi$}}}}
\fontsize{10}{0}\selectfont\put(177.165,33.7226){\makebox(0,0)[t]{\textcolor[rgb]{0.15,0.15,0.15}{{$-2 \pi$}}}}
\fontsize{10}{0}\selectfont\put(209.055,33.7226){\makebox(0,0)[t]{\textcolor[rgb]{0.15,0.15,0.15}{{$-\pi$}}}}
\fontsize{10}{0}\selectfont\put(272.835,33.7226){\makebox(0,0)[t]{\textcolor[rgb]{0.15,0.15,0.15}{{$\pi$}}}}
\fontsize{10}{0}\selectfont\put(304.724,33.7226){\makebox(0,0)[t]{\textcolor[rgb]{0.15,0.15,0.15}{{$2 \pi$}}}}
\fontsize{10}{0}\selectfont\put(336.614,33.7226){\makebox(0,0)[t]{\textcolor[rgb]{0.15,0.15,0.15}{{$3 \pi$}}}}
\fontsize{10}{0}\selectfont\put(368.504,33.7226){\makebox(0,0)[t]{\textcolor[rgb]{0.15,0.15,0.15}{{$4 \pi$}}}}
\fontsize{10}{0}\selectfont\put(400.394,33.7226){\makebox(0,0)[t]{\textcolor[rgb]{0.15,0.15,0.15}{{$5 \pi$}}}}
\fontsize{10}{0}\selectfont\put(432.283,33.7226){\makebox(0,0)[t]{\textcolor[rgb]{0.15,0.15,0.15}{{$6 \pi$}}}}
\fontsize{10}{0}\selectfont\put(235.948,150.562){\makebox(0,0)[r]{\textcolor[rgb]{0.15,0.15,0.15}{{1}}}}
\fontsize{10}{0}\selectfont\put(256.171,150.562){\makebox(0,0)[l]{\textcolor[rgb]{0,0,0}{{$\sinc(x)$}}}}
\fontsize{10}{0}\selectfont\put(454.606,33.5726){\makebox(0,0)[l]{\textcolor[rgb]{0,0,0}{{$x$}}}}
\end{picture}

  \caption{Gráfico da função $\sinc(x)$}
  \label{fig:functions_sinc}
\end{figure}
%------------------------------------------------------------------------------%

Derivada da sinc
\[
  \frac{d}{dx} \sinc(x) = \frac{\cos(x)-\sinc(x)}{x}
\]
Série de Taylor da sinc
\[
  \sinc(x) = \sum_{n=0}^{\infty} \frac{(-1)^{n} x^{2n}}{(2n+1)!}
\]
Cuidado para não confundir a definição utilizada aqui com
uma definição alternativa comumente utilizada no processamento digital de sinais,
da função \emph{sinc normalizada}\index{sinc!normalizada}
\[
  \sinc(x) = \frac{\sin(\pi x)}{\pi x}
\]

%------------------------------------------------------------------------------%
\subsection*{Outras Funções}

%------------------------------------------------------------------------------%
\begin{figure}
  \centering
  %--------------------------------------------------------------%
  \begin{subfigure}[b]{0.5\linewidth}
    \centering
    \input{fig_functions_carac}
    \caption{Gráfico da função característica}
    \label{fig:functions_carac}
  \end{subfigure}%
  \begin{subfigure}[b]{0.5\linewidth}
    \centering
    % Title: gl2ps_renderer figure
% Creator: GL2PS 1.4.2, (C) 1999-2020 C. Geuzaine
% For: Octave
% CreationDate: Fri Sep 10 20:09:34 2021
\setlength{\unitlength}{1pt}
\begin{picture}(0,0)
\includegraphics[width=8cm,height=4cm]{fig_functions_heaviside}
\end{picture}%
\begin{picture}(226,113)(0,0)
\fontsize{10}{0}\selectfont\put(17.7165,5.62978){\makebox(0,0)[t]{\textcolor[rgb]{0.15,0.15,0.15}{{-3}}}}
\fontsize{10}{0}\selectfont\put(49.6063,5.62978){\makebox(0,0)[t]{\textcolor[rgb]{0.15,0.15,0.15}{{-2}}}}
\fontsize{10}{0}\selectfont\put(81.4961,5.62978){\makebox(0,0)[t]{\textcolor[rgb]{0.15,0.15,0.15}{{-1}}}}
\fontsize{10}{0}\selectfont\put(145.276,5.62978){\makebox(0,0)[t]{\textcolor[rgb]{0.15,0.15,0.15}{{1}}}}
\fontsize{10}{0}\selectfont\put(177.165,5.62978){\makebox(0,0)[t]{\textcolor[rgb]{0.15,0.15,0.15}{{2}}}}
\fontsize{10}{0}\selectfont\put(209.055,5.62978){\makebox(0,0)[t]{\textcolor[rgb]{0.15,0.15,0.15}{{3}}}}
\fontsize{10}{0}\selectfont\put(108.384,91.6311){\makebox(0,0)[r]{\textcolor[rgb]{0.15,0.15,0.15}{{1}}}}
\fontsize{10}{0}\selectfont\put(119.764,103.406){\makebox(0,0)[l]{\textcolor[rgb]{0,0,0}{{$\heaviside(x)$}}}}
\fontsize{10}{0}\selectfont\put(221.811,7.6384){\makebox(0,0)[r]{\textcolor[rgb]{0,0,0}{{$x$}}}}
\end{picture}

    \caption{Gráfico da função de Heaviside}
    \label{fig:functions_heaviside}
  \end{subfigure}\\[0.5\baselineskip]
  %--------------------------------------------------------------%
  \begin{subfigure}[b]{0.5\linewidth}
    \centering
    \input{fig_functions_rampa}
    \caption{Gráfico da função rampa}
    \label{fig:functions_rampa}
  \end{subfigure}%
  \begin{subfigure}[b]{0.5\linewidth}
    \centering
    \input{fig_functions_delta}
    \caption{Representação da Delta de Dirac}
    \label{fig:functions_delta}
  \end{subfigure}\\[0.5\baselineskip]
  %--------------------------------------------------------------%
  \caption{Gráficos das funções característica, Heaviside, rampa e Delta de Dirac}
  \label{fig:functions_other}
\end{figure}
%------------------------------------------------------------------------------%

Definimos a \emph{Função Característica}\index{Função!característica}
de um subconjunto $A$ dos números reais $\R$ como a função $\carac{A}$ que vale 1
em $A$ e zero fora dele
\begin{equation}
  \label{eq:funcao_caracteristica}
  \carac{A}(x) = \begin{cases}
    \; 1 & \; x \in A \\
    \; 0 & \; \text{caso contrário}
  \end{cases}
\end{equation}

A \emph{Função de Heaviside}\index{Função!Heaviside}
ou \emph{Degrau Unitário}\index{Função!degrau unitário}
$\heaviside(x)$, vale 1 para $x$ positivo, zero para $x$ negativo.
\begin{equation}
  \label{eq:funcao_heaviside}
  \heaviside(x) = \begin{cases}
    \; 0 & \; x < 0 \\
    \; 1 & \; x > 0
  \end{cases}
\end{equation}
Essa função não precisa ser definida em $x=0$, mas dependendo do contexto
ela pode assumir algum valor específico nesse ponto. Um exemplo é
escolher o valor $\sfrac{1}{2}$
\[
  \heaviside(x) = \begin{cases}
    \;\; 0          & \; x < 0 \\
    \; \sfrac{1}{2} & \; x = 0 \\
    \;\; 1          & \; x > 0
  \end{cases}
\]
Também podemos escrever essa função como \(\heaviside(t) = \carac{[0, \infty)}(t)\).

A \emph{Função Rampa}\index{Função!rampa} é uma aproximação contínua para a
função de Heaviside, que é definida por
\begin{equation}
  \label{eq:funcao_rampa}
  \ramp_\varepsilon = \begin{cases}
    \;\; 0                                 & \; x < - \varepsilon \\[3mm]
    \; \dfrac{t}{2\varepsilon} + \dfrac{1}{2} & \; x = 0          \\[4mm]
    \;\; 1                                 & \; x > \varepsilon
  \end{cases}
\end{equation}

A \emph{Delta de Dirac}\index{Delta de Dirac}, \(\delta(t)\),
que representa um \emph{Impulso Unitário Instantâneo}\index{Impulso unitário instantâneo}
em \(t=0\). Ela é definida pela propriedade
\begin{equation}
  \label{eq:delta_dirac}
  \int_{-\infty}^\infty \delta(t)\; f(t)\, dt \;=\; f(0)
\end{equation}
Cuidado com o delta pois ele \emph{não é uma função} segundo a definição empregada no Cálculo.
Esclarecer seu significado preciso demanda a teoria de distribuições.
Mesmo assim, baseados em uma teoria matemática mais elaborada, escrevemos
informalmente que
\[
  \delta(t) = \lim_{h\to 0^+} d_h(t) =
    \begin{cases}
      \;\, 0      & \quad t\neq 0 \\
      \;   \infty & \quad t = 0
    \end{cases}
\]
A Seção~\ref{sec:convergencia_transformada_fourier} discute um pouco mais sobre essa função.

A Figura~\ref{fig:functions_other} mostra o gráfico dessas funções.
Porém, não podemos dizer que a Figura~\ref{fig:functions_delta} seja o gráfico
do Delta de Dirac. Essa é uma representação onde a seta indica que a função não
assume um valor real nesse ponto. Normalmente o valor $y$ onde a seta aponta
indica a constante multiplicando o delta.

%------------------------------------------------------------------------------%
